% QUESTAO
Use o teorema de Pitágoras para calcular a altura do paralelogramo, em seguida calcule a sua área:

\parbox{0.35\linewidth}{

% ALTERNATIVAS
$28$\,m$^2$
$30$\,m$^2$
$48$\,m$^2$
$24$\,m$^2$
$20$\,m$^2$

}\parbox{0.64\linewidth}{
\begin{center}
\vspace{-6mm}
\begin{tikzpicture}[ >=latex, scale=0.8 ]

\paralelogramo{}{7}{}{4}{3};
\coordinate (E) at ($(D)+(0,-4)$);
\draw[dashed] (D) -- (E);

\draw[nodelabel={-8pt}{$3$\,m}] (A) -- (E);
\draw[nodelabel={-8pt}{$4$\,m}] (E) -- (B);
\draw[nodelabel={-10pt}{$5$\,m}] (B) -- (C);

\draw (E) rectangle +(1*\raio,1*\raio);
\fill ($(E)+(0.5*\raio,0.5*\raio)$) circle (1*\raio mm);

\end{tikzpicture}
\vspace{-9mm}
\end{center}
}


